% This file was converted to LaTeX by Writer2LaTeX ver. 1.6.1
% see http://writer2latex.sourceforge.net for more info
\documentclass[a4paper]{article}
\usepackage[latin1]{inputenc}
\usepackage{amsmath}
\usepackage{amssymb,amsfonts,textcomp}
\usepackage[T1]{fontenc}
\usepackage[english]{babel}
\usepackage{color}
\usepackage{array}
\usepackage{hhline}
\usepackage{hyperref}
\hypersetup{pdftex, colorlinks=true, linkcolor=blue, citecolor=blue, filecolor=blue, urlcolor=blue, pdftitle=, pdfauthor=Ulrich Eisenecker, pdfsubject=, pdfkeywords=}
% Page layout (geometry)
\setlength\voffset{-1in}
\setlength\hoffset{-1in}
\setlength\topmargin{2cm}
\setlength\oddsidemargin{2cm}
\setlength\textheight{25.7cm}
\setlength\textwidth{17.001cm}
\setlength\footskip{0.0cm}
\setlength\headheight{0cm}
\setlength\headsep{0cm}
% Footnote rule
\setlength{\skip\footins}{0.119cm}
\renewcommand\footnoterule{\vspace*{-0.018cm}\setlength\leftskip{0pt}\setlength\rightskip{0pt plus 1fil}\noindent\textcolor{black}{\rule{0.25\columnwidth}{0.018cm}}\vspace*{0.101cm}}
% Pages styles
\makeatletter
\newcommand\ps@Standard{
  \renewcommand\@oddhead{}
  \renewcommand\@evenhead{}
  \renewcommand\@oddfoot{}
  \renewcommand\@evenfoot{}
  \renewcommand\thepage{\arabic{page}}
}
\makeatother
\pagestyle{Standard}
\title{}
\author{Ulrich Eisenecker}
\date{2024-06-03}
\begin{document}
This document describes the DSL for the specification of an SPL.

Version 1.1, June 3, 2024, by Yannis H�benthal \& Ulrich Eisenecker


\bigskip

Token types are comment, keywords, name and value.


\bigskip

A comment begins with // and must be followed by a whitespace (space, tab). It extends to the next newline character.


\bigskip

Keywords are SPL, FEATURES, CATEGORIES, END, MAX\_RANK, EXPLICIT\_DEPENDENCIES, MISSING\_SYSTEMS and AVAILABLE\_SYSTEMS.
A keyword must be terminated with a whitespace. Keywords are NOT case sensitive.

Strictly speaking, keywords are not reserved words. Nevertheless, it should be refrained from using them as names, e.g.
of features.


\bigskip

SPL has the cardinality 0..1. It enables the specification of a name for a DSL. The name must not contain any
whitespaces. A name is terminated by a whitespace.


\bigskip

FEATURES has the cardinality 1. It specifies the number of independent features of an SPL, followed by the names of the
independent features. There must be exactly as many feature names as the number of features specified. A number is
terminated by any character that cannot be part of an integer, e.g. a whitespace. IMPORTANT: The list of feature names
is NOT terminated with END. Currently, there is no check for duplicate feature names.


\bigskip

CATEGORIES has the cardinality 0..1. CATEGORIES is used to define which feature categories are to be included in the
feature analysis and feature isolation. Valid category values are F, A, O, AN, ON, ANF and ONF. If CATEGORY is not
explicitly specified, F is implicitly assumed to be the only category. The category values can be specified in any
order. If a category value is specified more than once, CATEGORY can appear anywhere in the SPL specification.


\bigskip

MAX\_RANK has the cardinality 0..1. It specifies the maximum number of feature interactions in relation to the
categories A, O, AN, ON, ANF and ONF that are taken into account. MAX\_RANK must be equal to or less than the number of
independent features. FEATURES must be specified before MAX\_RANK.


\bigskip

EXPLICIT\_DEPENDENCIES has the cardinality 0..1. It is used to specify explicitly dependent features. Dependencies are
REQUIRES\_ALL, REQUIRES\_ONE, REQUIRES\_ANY and EXCLUDES\_ALL. In addition, there is the special feature named ROOT
which represents the root feature that is a conceptual part of every system, but neither specified, nor analyzed or
isolated.

Each dependency is specified in the form feature\_name -- dependency -- list\_of\_feature\_names. The list of feature
names starts with at least one feature name and is terminated by END.

Examples:

ROOT REQUIRES\_ALL F1 F2 F3 END specifies F1, F2 and F3 as core features which occur in each system of the SPL.

F4 REQUIRES\_ALL F1 END specifies that F4 also requires F1 to be present in the system.

F5 REQUIRES\_ONE F6 F7 END specifies that F5 requires either F6 or F7 to be also included, which corresponds to an
exclusive-or group.

F8 REQUIRES\_ANY F9 F10 END specifies that F8 requires F9 or F10 or both, F9 and F10 to be included as well. This
corresponds to an inclusive-or group

F11 EXCLUDES\_ALL F12 F13 END specifies that the presence of F11 excludes the presence of F12 or F13.

The following specifications are syntactically correct and semantically valid, but redundant:

FA REQUIRES\_ALL FB END

FA REQUIRES\_ONE FB END

FA REQUIRES\_ANY FB END.

ROOT REQUIRES\_ONE FC FD END specifies that either FC or FD must be present in a system; its is not possible that both
are absent or present in a system.

After the parsing of the SPL specification is completed, a SAT solver is executed to check the specified dependencies
for consistency, also with regard to any specified available or missing systems. The list of explicit dependencies must
be concluded with END. If EXPLICIT\_DEPENDENCIES is missing, there are no explicitly dependent features.


\bigskip

MISSING\_SYSTEMS has a cardinality of 0..1 and specifies all systems that are not available. It must occur after
FEATURES. Each missing system is specified by its ID. The ID of the first system is 1, the ID of the last system is 2
to the power of F. In the system with ID 1 every feature is not available, in the system with the highest ID every
feature is available. The list of system IDs must be terminated with END. If EXPLICIT\_DEPENDENCIES is also specified,
no system ID may be specified that does not exist due to the constraints. It is permissible to specify missing systems
that would exist if the constraints were applied. The specification of MISSING\_SYSTEMS can be omitted. If
MISSING\_SYSTEMS occurs, AVAILABLE\_SYSTEMS may no longer be specified and vice versa. If neither MISSING\_SYSTEMS nor
AVAILABLE\_SYSTEMS are specified, it is assumed that all systems that can occur according to the constraints are
present.


\bigskip

AVAILABLE\_SYSTEMS has a cardinality of 0..1 and specifies all systems that are available. It must occur after FEATURES.
Each available system is specified by its ID. The ID of the first system is 1, the ID of the last system is 2 to the
power of F. In the system with ID 1 every feature is not available, in the system with the highest ID every feature is
available. The list of system IDs must be terminated with END. If EXPLICIT\_DEPENDENCIES is also specified, no system
ID may be specified that does not exist due to the constraints. It is permissible not to specify existing systems that
would exist if the constraints were applied. The specification of AVAILABLE\_SYSTEMS can be omitted. If
AVAILABLE\_SYSTEMS occurs, MISSING\_SYSTEMS may no longer be specified and vice versa. If neither MISSING\_SYSTEMS nor
AVAILABLE\_SYSTEMS are specified, it is assumed that all systems that can occur according to the constraints are
present.
\end{document}
